\documentclass[english,hebrew]{article}
\usepackage[HE8,T1]{fontenc}
\usepackage[utf8]{inputenc}
\setcounter{secnumdepth}{-2}
\setlength{\parindent}{0bp}
\usepackage{xcolor}
\usepackage{calc}
\usepackage{amsmath}
\usepackage{amsthm}
\usepackage{amssymb}

\makeatletter
\@ifundefined{date}{}{\date{}}
\usepackage{graphicx}
\usepackage{geometry}
\newcommand{\limi}{\underset{n\rightarrow\infty}{\lim}}
\newcommand{\limxz}{\underset{x\rightarrow x_{0}}{\lim}}
\newcommand{\limfxz}{\underset{x\rightarrow x_{0}}{\lim}f(x)}
\newcommand{\limfxm}{\underset{x\rightarrow x_{0}^{-}}{\lim}f(x)}
\newcommand{\limfxp}{\underset{x\rightarrow x_{0}^{+}}{\lim}f(x)}
\newcommand{\limxm}{\underset{x\rightarrow x_{0}^{-}}{\lim}}
\newcommand{\limxp}{\underset{x\rightarrow x_{0}^{+}}{\lim}}
\newcommand{\sqrm}{M_{m\times m}(\mathbb{F})}
\newcommand{\seqa}{(a_{n})_{n=1}^{\infty}}
\newcommand{\seqx}{(x_{n})_{n=1}^{\infty}}
\newcommand{\funcdr}{f:D\rightarrow\mathbb{R}}
\newcommand{\der}{\limxz\frac{f(x)-f(x_{0})}{x-x_{0}}}
\usepackage[all]{pdfprivacy}

\makeatother

\usepackage{babel}
\begin{document}
\newgeometry{top=1cm, bottom=2cm, left=1cm, right=1cm}
\title{אלגברה ליניארית {\beginL 2\endL} - הרצאה {\beginL 8\endL}}
\author{נעם לב}
\date{\L{$04.11.2026$}}

\maketitle
\noindent\fbox{\begin{minipage}[t]{1\columnwidth - 2\fboxsep - 2\fboxrule}%
\begin{center}
האוניברסיטה העברית \textbar{} אלגברה ליניארית ){\beginL 2\endL}( -
\L{$80135$} \textbar{} סמסטר{\small{} }{\beginL 2027\endL}{\small{} }א'
)תשפ\textquotedblleft ז א'( \textbar{} מרצה: פרופסור יבגני סטרחוב
\textbar{} מתרגל: מר לזר גוטרמן
\par\end{center}%
\end{minipage}}
\begin{center}
הסיכום הוא לא רשמי ולא עבר אישור של הסגל, על כן אין להשתמש בו )ובין
השאר במספור שלו ובשמות הטענות שבו( כמקור רשמי. אם ישנן טעויות, הן
שלי ולא של הסגל.
\par\end{center}

\medskip{}

\tableofcontents{}

\newpage{}

\section{הקדמה - מה נשאר בנושא הלכסון}

אנחנו מדברים על מתי אפשר ללכסן העתקות לינאריות )וספציפית אופרטורים(.

מה שנשאר לעשות בנושא הזה הוא להוכיח שני משפטים: אחד על השימוש בפולינום
האופייני, והשני על השימוש בפולינום המינימלי.

לשם כך צריך להגדיר שני דברים: ריבוי אלגברי וריבוי גיאומטרי.\\

רלוונטי לכל ההרצאה )אותה הנחה( ולכן נציין רק כאן: \\

יהי \L{$V$} מ\textquotedblright ו מעל שדה \L{$\mathbb{F}$} כלשהו
כך ש-\L{%
\[
\dim V=n\geq1,\quad T\in\mathcal{L}(V)
\]
}%

\part{שימוש בפולינום האופייני ללכסון )ריבוי אלגברי וריבוי גיאומטרי(}

\section{הגדרה - ריבוי אלגברי של ערך עצמי \L{$(\text{Algebraic multiplicity})$}}

יהי \L{$\lambda$} ערך עצמי של \L{$T$}.

המספר \L{$m_{a}(\lambda)$} הגדול ביותר כך ש- \L{$(z-\lambda)^{m_{a}(\lambda)}$}
מחלק את הפולינום האופייני נקרא הריבוי האלגברי של \L{$\lambda$}.\\

\L{$\lambda$} הוא ערך עצמי, ולכן הוא שורש של הפולינום האופייני )הרצאה
\L{$2$}(, כלומר \L{$(z-\lambda)$} מחלק את הפולינום האופייני.

מכאן, תמיד מתקיים \L{$m_{a}(\lambda)\geq1$}.\\

יבגני: יש יותר מדי \L{$m$} בקורס הזה, אבל זה הסימון הסטנדרטי. מי
שהסימון לא נראה לו יכול להשתמש בסימון אחר. )-\_-(\\

נעם: מעתה אאמץ את הסימונים מהתרגול: \L{$a_{\lambda}$} עבור הריבוי
האלגברי של \L{$\lambda$} ו-\L{$g_{\lambda}$} עבור הריבוי הגיאומטרי
של \L{$\lambda$} )במקום \L{$m_{a}(\lambda)$} ו-\L{$m_{geom}(\lambda)$}
בהתאמה, מההרצאה(.

\section{דוגמה - ריבוי אלגברי מעל \L{$\mathbb{R}$}}

\L{%
\[
\mathbb{F}=\mathbb{R}
\]
}%
\L{%
\[
P_{ch}(z)=(z-2)^{3}(z+1)(z^{2}+2z+2)
\]
}%

לפולינום \L{$z^{2}+2z+2$} אין שורשים ממשיים )הדיסקרימיננטה היא \L{$4-8<0$}(,
ולכן מעל \L{$\mathbb{R}$} הערכים העצמיים היחידים הם \L{$2$} ו-\L{$\left(-1\right)$}.

ייתכן מצב כזה.

אז

\L{%
\[
a_{\left(2\right)}=3,\quad a_{\left(-1\right)}=1
\]
}%

כלומר, הריבוי האלגברי הוא החזקה שמתקבלת בפולינום האופייני.

\section{הגדרה - ריבוי גאומטרי של ערך עצמי}

יהי \L{$\lambda$} ערך עצמי של \L{$T$}.

אז המימד של התמ\textquotedblright ו שמשרה \L{$\lambda$} הוא הריבוי
הגיאומטרי של \L{$\lambda$}:\L{%
\[
g_{\left(\lambda\right)}\overset{\text{def}}{=}\dim\left(\text{Null}(T-\lambda I)\right)
\]
}%

למה צריך את זה? מתברר שיש משפט שימושי מאוד, אולי השימושי ביותר בכל
התחום הזה.

\section{משפט - לכסינות באמצעות ריבוי אלגברי וריבוי גאומטרי}

יהיו \L{$\lambda_{1},\dots,\lambda_{m}$} כל הערכים העצמיים של \L{$T$}
)שונים זה מזה(, כלומר\L{%
\[
\forall i,j\in\{1,\dots,m\}\text{ s.t. }i\neq j:\quad\lambda_{j}\neq\lambda_{i}
\]
}%

אז התנאים הבאים שקולים:

)I( קיים בסיס \L{$\mathcal{B}_{V}$} ב-\L{$V$} כך ש- \L{%
\[
M(T,\mathcal{B}_{V},\mathcal{B}_{V})=\begin{pmatrix}a_{1,1} & 0_{\mathbb{F}} & \cdots & 0_{\mathbb{F}}\\
0_{\mathbb{F}} & a_{2,2} &  & \vdots\\
\vdots &  & \ddots\\
0_{\mathbb{F}} & 0_{\mathbb{F}} &  & a_{n,n}
\end{pmatrix}
\]
}%

)II( מתקיימים שני התנאים הבאים:

)II( )א( סכום הריבויים האלגבריים של הערכים העצמיים שווה למימד של \L{$V$}:\L{%
\[
a_{\left(\lambda_{1}\right)}+\dots+a_{\left(\lambda_{m}\right)}=n
\]
}%

)II( )ב( עבור כל ערך עצמי, הריבוי האלגברי שלו שווה לריבוי הגיאומטרי
שלו:

\L{%
\[
g_{\left(\lambda_{1}\right)}=a_{\left(\lambda_{1}\right)},\dots,g_{\left(\lambda_{m}\right)}=a_{\left(\lambda_{m}\right)}
\]
}%

כלומר, התנאי השני הוא )א( ו-)ב( ביחד.
\begin{proof}
הכיוון \L{$(\text{I})\Rightarrow(\text{II})$}:

1. לפי המשפט על התנאים השקולים ללכסינות מהרצאה {\beginL 7\endL}, קיים
בסיס \L{$\mathcal{B}_{V}$} המורכב מוקטורים עצמיים של \L{$T$}.

2. נסדר את הוקטורים מחדש לפי הערכים העצמיים שלהם )כמו בהרצאה {\beginL 7\endL}(.

לכל \L{$i\in\{1,\dots,m\}$}, נסמן את וקטורי הבסיס השייכים ל- \L{$\lambda_{i}$}
כך: \L{%
\[
\mathcal{B}^{(i)}=(v_{1}^{(i)},\dots,v_{n_{i}}^{(i)})
\]
}%
האינדקס העליון אומר לאיזה ערך עצמי הוקטור שייך, והתחתון ממספר את הוקטורים
בתוך הקבוצה. 

נשים את הקבוצות אחת אחרי השנייה: \L{%
\[
\mathcal{B}'_{V}=(\mathcal{B}^{(1)},\mathcal{B}^{(2)},\dots,\mathcal{B}^{(m)})
\]
}%
\L{$\mathcal{B}'_{V}$} הוא בסיס של \L{$V$}, כי אלה אותם וקטורים
בסדר אחר.

לא החסרנו ולא הוספנו וקטורים מהבסיס לכן\L{%
\[
n_{1}+\dots+n_{m}=n
\]
}%

3. מאיך שהגדרנו אותו, לכל וקטור בבסיס \L{$v_{j}^{(i)}$} מתקיים \L{$T\left(v_{j}^{(i)}\right)=\lambda_{i}v_{j}^{(i)}$}.

אז מהגדרת המטריצה המייצגת \L{$\lambda_{1}$} מופיע \L{$n_{1}$} פעמים,
\L{$\lambda_{2}$} מופיע \L{$n_{2}$} פעמים וכן הלאה עד \L{$\lambda_{m}$}
שמופיע \L{$n_{m}$} פעמים, ובכל מקום אחר יש אפס.

\L{%
\[
M(T,\mathcal{B}'_{V},\mathcal{B}'_{V})=\begin{pmatrix}\lambda_{1} &  &  &  &  &  & 0_{\mathbb{F}}\\
 & \ddots\\
 &  & \lambda_{1}\\
 &  &  & \ddots\\
 &  &  &  & \lambda_{m}\\
 &  &  &  &  & \ddots\\
0_{\mathbb{F}} &  &  &  &  &  & \lambda_{m}
\end{pmatrix}
\]
}%

מכאן, \L{%
\[
P_{ch}(z)=\det\left(zI_{n}-M(T,\mathcal{B}'_{V},\mathcal{B}'_{V})\right)=(z-\lambda_{1})^{n_{1}}\cdots(z-\lambda_{m})^{n_{m}}
\]
}%
\L{$\lambda_{1},\dots,\lambda_{m}$} שונים זה מזה, לכן החזקה הגדולה
ביותר של \L{$(z-\lambda_{i})$} שמחלקת את \L{$P_{ch}$} היא בדיוק
\L{$n_{i}$} ונקבל

\L{%
\[
a_{\left(\lambda_{1}\right)}=n_{1},\dots,a_{\left(\lambda_{m}\right)}=n_{m}
\]
}%
\L{%
\[
\therefore\quad a_{\left(\lambda_{1}\right)}+\dots+a_{\left(\lambda_{m}\right)}=n_{1}+\dots+n_{m}=n
\]
}%

שזה )א(.

{\beginL 4\endL}. כל \L{$v_{j}^{(i)}$} הוא וקטור עצמי של \L{$\lambda_{i}$},
כלומר שייך ל- \L{$\text{Null}(T-\lambda_{i}I)$}, ו- \L{$\text{Null}(T-\lambda_{i}I)$}
הוא תמ\textquotedblright ו ולכן מכיל גם את כל הצירופים הלינאריים שלהם.\L{%
\begin{align*}
\text{Span}\left(v_{1}^{(1)},\dots,v_{n_{1}}^{(1)}\right) & \subseteq\text{Null}(T-\lambda_{1}I)\\
 & \vdots\\
\text{Span}\left(v_{1}^{(m)},\dots,v_{n_{m}}^{(m)}\right) & \subseteq\text{Null}(T-\lambda_{m}I)
\end{align*}
}%

{\beginL 5\endL}. \L{$\mathcal{B}'_{V}$} הוא בסיס של \L{$V$} שמחולק
לקבוצות \L{$\mathcal{B}^{(1)},\dots,\mathcal{B}^{(m)}$}, להן אין
וקטורים משותפים.

כל \L{$v\in V$} נכתב באופן יחיד כצירוף לינארי של וקטורי הבסיס.

אם מקבצים את המחוברים לפי הקבוצות מקבלים הצגה יחידה \L{$v=u_{1}+\dots+u_{m}$}
כאשר \L{$u_{i}\in\text{Span}\left(\mathcal{B}^{(i)}\right)$}.

לכן \L{$V$} הוא סכום ישר של ה-\L{$\text{Span}$}-ים )זו בדיוק ההגדרה
של אקסלר לסכום ישר(.

\L{%
\[
V=\text{Span}\left(v_{1}^{(1)},\dots,v_{n_{1}}^{(1)}\right)\oplus\dots\oplus\text{Span}\left(v_{1}^{(m)},\dots,v_{n_{m}}^{(m)}\right)
\]
}%

{\beginL 6\endL}. לפי המשפט על התנאים השקולים ללכסינות מהרצאה \L{$7$},
\L{%
\[
V=\text{Null}(T-\lambda_{1}I)\oplus\dots\oplus\text{Null}(T-\lambda_{m}I)
\]
}%

7. נרצה להוכיח שלכל \L{$i\in\{1,\dots,m\}$} מתקיים \L{%
\[
\text{Span}\left(v_{1}^{(i)},\dots,v_{n_{i}}^{(i)}\right)=\text{Null}(T-\lambda_{i}I)
\]
}%
נניח בשלילה שזה לא נכון, כלומר שקיים \L{$i\in\{1,\dots,m\}$} כך ש-
\L{%
\[
\text{Span}\left(v_{1}^{(i)},\dots,v_{n_{i}}^{(i)}\right)\neq\text{Null}(T-\lambda_{i}I)
\]
}%
לפי ){\beginL 4\endL}( מתקיים \L{$\text{Span}\left(\mathcal{B}^{(i)}\right)\subseteq\text{Null}(T-\lambda_{i}I)$},
לכן בהכרח \L{$\text{Null}(T-\lambda_{i}I)\not\subseteq\text{Span}\left(\mathcal{B}^{(i)}\right)$}.

אז קיים \L{$u\in V$} כך ש- \L{%
\[
u\in\text{Null}(T-\lambda_{i}I)
\]
}%
אבל \L{%
\[
u\notin\text{Span}\left(\mathcal{B}^{(i)}\right)
\]
}%
בפרט, \L{$0_{V}\in\text{Span}\left(\mathcal{B}^{(i)}\right)$} ולכן
\L{$u\neq0_{V}$}.

לפי ){\beginL 5\endL}(, קיימים \L{$u_{1},\dots,u_{m}$} כך ש-

\L{%
\[
u=\underset{\in\text{Span}\left(v_{1}^{(1)},\dots,v_{n_{1}}^{(1)}\right)}{\underbrace{u_{1}}}+\quad\dots\quad+\underset{\in\text{Span}\left(v_{1}^{(m)},\dots,v_{n_{m}}^{(m)}\right)}{\underbrace{u_{m}}}
\]
}%

קיים \L{$i\ne\ell\in\{1,\dots,m\}$}, כך ש- \L{$u_{\ell}\neq0_{V}$},
אחרת כולם וקטור האפס למעט \L{$u_{i}$} ונקבל ש-\L{$u=u_{i}\in\text{Span}\left(\mathcal{B}^{(i)}\right)$}
וזאת סתירה.

נקבל \L{%
\[
0_{V}\neq u_{\ell}\in\text{Span}\left(v_{1}^{(\ell)},\dots,v_{n_{\ell}}^{(\ell)}\right)\underset{(4)}{\subseteq}\text{Null}(T-\lambda_{\ell}I)
\]
}%
נעביר את \L{$u$} אגף.

לפי ){\beginL 4\endL}(, לכל \L{$k\in\left\{ 1,\dots,m\right\} $}
מתקיים \L{$u_{k}\in\text{Null}(T-\lambda_{k}I)$}.

בנוסף, \L{$u\in\text{Null}(T-\lambda_{i}I)$}.

בסה\textquotedblright כ נקבל\L{%
\[
0_{V}=\underset{\in\text{Null}(T-\lambda_{1}I)}{\underbrace{u_{1}}}+\dots+\underset{\in\text{Null}(T-\lambda_{i}I)}{\underbrace{(u_{i}-u)}}+\dots+\underset{\in\text{Null}(T-\lambda_{m}I)}{\underbrace{u_{m}}}
\]
}%

נזכיר ש- \L{$u_{\ell}\neq0_{V}$}.

מצד שני,\L{%
\[
0_{V}=\underset{\in\text{Null}(T-\lambda_{1}I)}{\underbrace{0_{V}}}+\dots+\underset{\in\text{Null}(T-\lambda_{m}I)}{\underbrace{0_{V}}}
\]
}%
קיבלנו שתי הצגות שונות של \L{$0_{V}$} כסכום של וקטורים מ- \L{$\text{Null}(T-\lambda_{1}I),\dots,\text{Null}(T-\lambda_{m}I)$}.

זאת סתירה ל-){\beginL 6\endL}(, כי בסכום ישר ההצגה יחידה.

לכן ההנחה לא נכונה וקיבלנו: \L{%
\begin{align*}
\text{Span}\left(v_{1}^{(1)},\dots,v_{n_{1}}^{(1)}\right) & =\text{Null}(T-\lambda_{1}I)\\
 & \vdots\\
\text{Span}\left(v_{1}^{(m)},\dots,v_{n_{m}}^{(m)}\right) & =\text{Null}(T-\lambda_{m}I)
\end{align*}
}%
{\beginL 8\endL}. הוקטורים של \L{$\mathcal{B}^{(i)}$} הם חלק מהבסיס
\L{$\mathcal{B}'_{V}$} ולכן בת\textquotedblright ל.

לכן הרשימה \L{$\mathcal{B}^{(i)}=\left(v_{1}^{(i)},\dots,v_{n_{i}}^{(i)}\right)$}
היא בת\textquotedblright ל ופורשת, דהיינו בסיס של \L{$\text{Span}\left(\mathcal{B}^{(i)}\right)$},
ומתקיים\L{%
\[
\dim\left(\text{Span}\left(\mathcal{B}^{(i)}\right)\right)=n_{i}
\]
}%

לפי ){\beginL 7\endL}(, מתקיים \L{$\dim\left(\text{Null}(T-\lambda_{i}I)\right)=\dim\left(\text{Span}\left(\mathcal{B}^{(i)}\right)\right)=n_{i}$}

9. לפי ){\beginL 3\endL}( מתקיים \L{$a_{(\lambda_{i})}=n_{i}$}, ולפי
ההגדרה של הריבוי הגאומטרי, \L{$g_{(\lambda_{i})}=\dim\left(\text{Null}(T-\lambda_{i}I)\right)$}.

אז בסה\textquotedblright כ\L{%
\[
\forall1\le i\le m:
\]
}%
\L{%
\[
a_{(\lambda_{i})}=n_{i}=\dim\left(\text{Null}(T-\lambda_{i}I)\right)=g_{(\lambda_{i})}
\]
}%
שזה )ב(, כנדרש.

יבגני, בתשובה לשאלות בהפסקה: ההוכחה הזו יכולה להופיע במבחן.

לא מספיק להגיע עד שלבים 4, 5, 6 ולכתוב \textquotedblright ולכן יש
שוויון\textquotedblright ; על זה תרד נקודה, כי זה לא ברור מאליו.

צריך לכתוב לפחות כמה מילים: שאם בוחרים וקטור מכאן אז מגיעים לוקטור
שם, וזה לא יכול להיות.

ההוכחה של יבגני בכיתה הייתה הוכחה קצת שונה מזו שכתבתי כאן.

ניתן להוכיח גם ישירות )נמצא בנספח(.
\end{proof}
\begin{proof}
הכיוון \L{$(\text{II})\Rightarrow(\text{I})$}:

1. מ-)ב(: \L{%
\begin{align*}
a_{\left(\lambda_{1}\right)}= & \dim\text{Null}(T-\lambda_{1}I)\\
 & \vdots\\
a_{\left(\lambda_{m}\right)}= & \dim\text{Null}(T-\lambda_{m}I)
\end{align*}
}%
ולכן מ-)א(: \L{%
\[
\dim\text{Null}(T-\lambda_{1}I)+\dots+\dim\text{Null}(T-\lambda_{m}I)=\dim V
\]
}%

זה תנאי \L{$\left(IV\right)$} במשפט על תנאים שקולים ללכסינות, לכן
\L{$T$} לכסין כנדרש.
\end{proof}
לפי יבגני, המעבר הזה הוא המטרה העיקרית של המשפט הקודם: התנאי די חלש,
מסתכלים רק על מימדים, ופתאום משוויון המימדים אפשר לומר שהמטריצה לכסינה.\\

יבגני: במבחן לא צריך להוכיח שוב את המשפט הקודם, אבל צריך לצטט אותו
במפורש {\beginL -\endL} \textquotedblright הייתה טענה שמשוויון המימדים
נובע שניתן להציג את \L{$T$} על ידי מטריצה אלכסונית\textquotedblright .
זה הכלל תמיד: מותר להסתמך על משפטים קודמים מהרשימה בלי הוכחה, אבל
צריך לצטט את הטענה שמתבססים עליה.\\

זהו השימוש בפולינום האופייני.

מחקירת הפולינום האופייני מקבלים מיד את הריבוי האלגברי של כל ערך עצמי
של \L{$T$}. 

אחר כך צריך לעשות מאמץ ולחשב גם את המימדים של תתי-המרחבים \L{$\text{Null}(T-\lambda_{i}I)$},
כלומר למצוא בהם בסיס של וקטורים עצמיים.

מי שמצליח בזה יכול לקבוע אם יש ייצוג על ידי מטריצה אלכסונית.

לפי יבגני, זה התנאי הכי פרקטי שיש לבדיקה אם ה\textquotedblright ל
מעל שדה כלשהו לכסינה.

יש גם תנאי שמשתמש בפולינום המינימלי, ונוכיח אותו אחרי ההפסקה.

דוגמאות נוספות יהיו כנראה בתרגול.

\section{דוגמה - אופרטור שאינו לכסין}

\L{%
\[
V=\mathbb{R}^{2}=\left\{ \begin{pmatrix}x\\
y
\end{pmatrix}\ \middle|\ x\in\mathbb{R},\ y\in\mathbb{R}\right\} 
\]
}%
\L{%
\[
T_{A}\in\mathcal{L}(\mathbb{R}^{2})
\]
}%
\L{%
\[
T_{A}(v)=\begin{pmatrix}1 & 1\\
0 & 1
\end{pmatrix}\cdot v
\]
}%

ההעתקה מוגדרת על ידי מטריצה, וברור שהיא לינארית. 

בשאלה כזו תתבקשו לחקור אם המטריצה לכסינה, כלומר אם ניתן להציג את ההעתקה
על ידי מטריצה אלכסונית.\\

טענה: \L{$T_{A}$} אינו לכסין.

יש כמה אפשרויות, אפשר גם לחשב כאן את הפולינום המינימלי, אבל נשתמש
במשפט הקודם.
\begin{proof}
1. \L{%
\[
P_{ch}(z)=\det\begin{pmatrix}z-1 & -1\\
0 & z-1
\end{pmatrix}=(z-1)^{2}
\]
}%

המטריצה משולשית עליונה, ולכן הדטרמיננטה היא מכפלת איברי האלכסון.

אזי \L{$\lambda=1$} הוא ערך עצמי של \L{$T_{A}$}, ו- \L{%
\[
a_{\left(1\right)}=2
\]
}%

2. כדי למצוא את הריבוי הגאומטרי צריך למצוא בסיס לגרעין, כלומר בסיס
שבנוי מוקטורים עצמיים של הערך העצמי שלנו. \L{%
\[
\begin{pmatrix}1 & 1\\
0 & 1
\end{pmatrix}\begin{pmatrix}x\\
y
\end{pmatrix}=1\cdot\begin{pmatrix}x\\
y
\end{pmatrix}
\]
}%
כלומר \L{%
\[
\begin{cases}
x+y=x\\
y=y
\end{cases}
\]
}%
ולכן \L{%
\[
y=0
\]
}%
\L{%
\[
\therefore\quad\text{Null}(T_{A}-1\cdot I)=\text{Span}\left(\begin{pmatrix}1\\
0
\end{pmatrix}\right)
\]
}%

\L{%
\[
\therefore\quad g_{\left(1\right)}=1
\]
}%

כלומר\L{%
\[
a_{\left(1\right)}\neq g_{\left(1\right)}
\]
}%
ולפי המשפט שלמדנו \L{$A$} אינה לכסינה.
\end{proof}
יבגני: זה התרגיל הטיפוסי, כשלא רוצים לתת לכם להסתבך עם חיפוש של וקטורים
עצמיים.

מבחינה טכנית, הדבר הקשה ביותר בו הוא לפתור מערכת משוואות.

\part{שימוש בפולינום המינימלי ללכסון}

\section{משפט - לכסינות והפולינום המינימלי}

יהי \L{$q_{\min}(z)$} הפולינום המינימלי של \L{$T$}.

אז התנאים הבאים שקולים:\\

)I( קיים בסיס \L{$\mathcal{B}_{V}$} כך ש- \L{%
\[
M(T,\mathcal{B}_{V},\mathcal{B}_{V})=\begin{pmatrix}a_{1,1} &  & 0_{\mathbb{F}}\\
 & \ddots\\
0_{\mathbb{F}} &  & a_{n,n}
\end{pmatrix}
\]
}%

)II( \L{%
\[
q_{\min}(z)=(z-\lambda_{1})\cdots(z-\lambda_{m})
\]
}%
כאשר \L{$\lambda_{1},\dots,\lambda_{m}$} הם ערכים עצמיים שונים.

כלומר, ניתן לכתוב את הפולינום המינימלי כמכפלה של פולינומים לינאריים,
כאשר \L{$\lambda_{1},\dots,\lambda_{m}$} שונים זה מזה )זה ההבדל מהמשפט
מהרצאה \L{$6$} על ייצוג משולשי(.
\begin{proof}
הכיוון \L{$(\text{I})\Rightarrow(\text{II})$}:

1. \L{$T$} לכסין ולכן לפי המשפט מהרצאה {\beginL 7\endL} על תנאים
שקולים ללכסינות יש ל-\L{$V$} בסיס \L{$\mathcal{B}_{V}=(v_{1},\dots,v_{n})$}
של וקטורים עצמיים של \L{$T$}.

נסמן ב- \L{$\lambda_{1},\dots,\lambda_{m}$} את כל הערכים העצמיים
השונים של \L{$T$}.

נרצה להוכיח שעבורם מתקיים \L{%
\[
q_{\min}(z)=(z-\lambda_{1})\cdots(z-\lambda_{m})
\]
}%

{\beginL 2\endL}.\L{%
\begin{align*}
\left((T-\lambda_{1}I)\cdots(T-\lambda_{m}I)\right)\left(v_{1}\right) & \underset{\left(*\right)}{=}0_{V}\\
 & \vdots\\
\left((T-\lambda_{1}I)\cdots(T-\lambda_{m}I)\right)\left(v_{n}\right) & \underset{\left(**\right)}{=}0_{V}
\end{align*}
}%

לא יודעים מהו הערך העצמי של \L{$v_{\ell}$}, אבל הוא אחד מה-\L{$\lambda_{i}$}:

\L{%
\[
\left(*\right)\quad\exists1\le k\le m:(T-\lambda_{k}I)\left(v_{1}\right)=0_{V}
\]
}%
\L{%
\[
\left(**\right)\quad\exists1\le k\le m:(T-\lambda_{k}I)\left(v_{n}\right)=0_{V}
\]
}%

הגורמים מתחלפים זה עם זה כי כולם פולינומים באותו אופרטור \L{$T$}
)הרצאה {\beginL 5\endL}: הערה - מדוע מותר להציב \L{$T$} במקום \L{$x$}(,
ולכן תמיד אפשר לסדר את המכפלה כך שהגורם המתאים \L{$(T-\lambda_{k}I)$}
יפעל ראשון על \L{$v_{\ell}$} ויאפס אותו, ואז כמובן ששאר הגורמים מחזירים
את וקטור האפס כאשר הם מופעלים על וקטור האפס.

אבל מכך ש-\L{$\mathcal{B}_{V}$} הוא בסיס, כל וקטור ב-\L{$V$} הוא
צירוף לינארי של הוקטורים האלה וקיבלנו ש-\L{%
\[
\forall v\in V:\quad\left((T-\lambda_{1}I)\cdots(T-\lambda_{m}I)\right)\left(v\right)=0_{V}
\]
}%

{\beginL 3\endL}. כלומר\L{%
\[
(T-\lambda_{1}I)\cdots(T-\lambda_{m}I)=0_{\mathcal{L}(V)}
\]
}%

{\beginL 4\endL}. ראינו בהרצאה חמש שהפולינום המינימלי מחלק כל פולינום
שמאפס את \L{$T$}, לכן הפולינום המינימלי בפרט מחלק את הפולינום הזה.

{\beginL 5\endL}. ראינו בהרצאה {\beginL 5\endL} )משפט - שורשי הפולינום
המינימלי הם הערכים העצמיים( שהערכים העצמיים של \L{$T$} מאפסים את
הפולינום המינימלי.

{\beginL 6\endL}. כלומר מצד אחד \L{$q_{\min}$} צריך לחלק את המכפלה,
ולכן הוא המכפלה של כל הגורמים או של חלק מהם.

מצד שני הוא חייב להתאפס בכל הנקודות \L{$\lambda_{1},\dots,\lambda_{m}$}.

לכן אין ברירה, והוא שווה למכפלה כולה.

בסך הכל קיבלנו \L{%
\[
q_{\min}(z)=(z-\lambda_{1})\cdots(z-\lambda_{m})
\]
}%
הכיוון \L{$(\text{II})\Rightarrow(\text{I})$}:

נוכיח באינדוקציה לפי \L{$m$}.

נרצה להוכיח שב-\L{$V$} יש בסיס המורכב מוקטורים עצמיים של \L{$T$},
שזה אותו דבר כמו לומר שהמטריצה אלכסונית.

1. נניח כי \L{$m=1$}. אזי \L{%
\[
q_{\min}(z)=z-\lambda_{1}
\]
}%
\L{%
\[
q_{\min}(T)=T-\lambda_{1}I=0
\]
}%
\L{%
\[
T=\lambda_{1}I
\]
}%
כלומר \L{%
\[
M(T,\mathcal{B}_{V},\mathcal{B}_{V})=\begin{pmatrix}\lambda_{1} &  & 0_{\mathbb{F}}\\
 & \ddots\\
0_{\mathbb{F}} &  & \lambda_{1}
\end{pmatrix}
\]
}%

כנדרש.

2. יהי \L{$m\geq2$}. נניח כי המשפט מתקיים לכל מ\textquotedblright ו
נ\textquotedblright ס \L{$\widetilde{V}$} מעל אותו שדה \L{$\mathbb{F}$}
ואופרטור \L{$\widetilde{T}\in\mathcal{L}(\widetilde{V})$}, כך ש-
\L{%
\[
\widetilde{q}_{\min}(z)=(z-\alpha_{1})\cdots(z-\alpha_{k})
\]
}%
כאשר \L{$\alpha_{1},\dots,\alpha_{k}$} שונים זה מזה ו- \L{%
\[
1\leq k\leq m-1
\]
}%
נרצה לבנות בסיס \L{$\mathcal{B}_{V}$} ב-\L{$V$} המורכב מוקטורים
עצמיים של \L{$T$}.

\textcolor{teal}{חשוב מאוד לכתוב את הנחת האינדוקציה נכון; היא מאוד
משמעותית כאן.}

\textcolor{teal}{הטענה שמוכיחים באינדוקציה היא: כל אופרטור, על כל
מרחב נ'’ס מעל \L{$\mathbb{F}$}, שהפולינום המינימלי שלו הוא מכפלה
של \L{$m$} גורמים לינאריים שונים, הוא לכסין.}

\textcolor{teal}{לא משנה מהו המרחב, מהי ההעתקה, מהו המימד ומהו הפולינום
המינימלי שלה.}

\textcolor{teal}{חשוב רק שאפשר להציג אותו כמכפלה כזו, ושמספר הגורמים
קטן יותר.}

\textcolor{teal}{נעבוד לא עם \L{$V$} המקורי אלא עם מרחב אחר, ועם
הצמצום של ההעתקה אליו, שהוא פורמלית העתקה אחרת.}

\textcolor{teal}{השדה נשאר אותו שדה.}

3. נגדיר \L{%
\[
U=\text{Range}(T-\lambda_{m}I)
\]
}%
נרצה להפעיל את הנחת האינדוקציה )שלב 2( על \L{$\widetilde{V}=U$} ו-\L{$\widetilde{T}=T|_{U}$}. 

לשם כך נוכיח שלושה דברים.

)א( \L{$U$} הוא תמ\textquotedblright ו שמור תחת \L{$T$}, ולכן \L{$T|_{U}\in\mathcal{L}(U)$}:

למעשה, הוכחנו זאת כבר כמה פעמים.

למשל, הראינו זאת בהוכחת המשפט על ייצוג משולשי מעל \L{$\mathbb{C}$}
בהרצאה \L{$3$}.

לכל \L{$u\in U$} מתקיים \L{%
\[
T(u)=(T-\lambda_{m}I)(u)+\lambda_{m}u\in U
\]
}%
כי \L{$(T-\lambda_{m}I)(u)\in\text{Range}(T-\lambda_{m}I)=U$}, ו-\L{$\lambda_{m}u\in U$}
כי \L{$U$} תמ\textquotedblright ו.

לכן הצמצום של \L{$T$} ל-\L{$U$} הוא אופרטור על \L{$U$}: \L{%
\[
T|_{U}\in\mathcal{L}(U)
\]
}%

)ב( \L{$\dim U\geq1$}:

נניח בשלילה ש- \L{$U=\{0_{V}\}$}.

אז לכל \L{$v\in V$} מתקיים \L{%
\[
(T-\lambda_{m}I)(v)=0_{V}
\]
}%
כלומר הפולינום \L{$z-\lambda_{m}$}, שמעלתו \L{$1$}, מאפס את \L{$T$}. 

זאת סתירה לכך שהפולינום המינימלי של \L{$T$} הוא ממעלה \L{$m\geq2$},
כי הוא הפולינום ממעלה מינימלית שמאפס את \L{$T$}.

לכן \L{$U\neq\{0_{V}\}$}, כלומר \L{%
\[
\dim U\geq1
\]
}%

)ג( הפולינום המינימלי של \L{$T|_{U}$} הוא מכפלה של \L{$k$} גורמים
לינאריים שונים זה מזה, כאשר \L{$1\leq k\leq m-1$}:

יהי \L{$u\in U$}. לפי הגדרת הטווח, קיים \L{$v_{0}\in V$} כך ש- \L{%
\[
u=(T-\lambda_{m}I)(v_{0})
\]
}%
נפעיל על \L{$u$} את \L{$m-1$} הגורמים הראשונים.

יחד עם הגורם \L{$(T-\lambda_{m}I)$} מקבלים את כל \L{$m$} הגורמים,
כלומר את \L{$q_{\min}(T)$} )מההנחה(: \L{%
\begin{align*}
\left((T-\lambda_{1}I)\cdots(T-\lambda_{m-1}I)\right)(u) & =\left((T-\lambda_{1}I)\cdots(T-\lambda_{m-1}I)(T-\lambda_{m}I)\right)(v_{0})\\
 & =q_{\min}(T)(v_{0})\\
 & =0_{V}
\end{align*}
}%
כי \L{$q_{\min}(z)=(z-\lambda_{1})\cdots(z-\lambda_{m})$} הוא הפולינום
המינימלי של \L{$T$}, ולכן \L{$q_{\min}(T)=0_{\mathcal{L}(V)}$}. 

קיבלנו \L{%
\[
\forall u\in U:\quad\left((T-\lambda_{1}I)\cdots(T-\lambda_{m-1}I)\right)(u)=0_{V}
\]
}%
נגדיר \L{%
\[
p(z)=(z-\lambda_{1})\cdots(z-\lambda_{m-1})
\]
}%
לכל \L{$u\in U$} מתקיים \L{$T|_{U}(u)=T(u)$}, ולכן לפי מה שקיבלנו
\L{%
\[
p\left(T|_{U}\right)(u)=\left((T-\lambda_{1}I)\cdots(T-\lambda_{m-1}I)\right)(u)=0_{V}
\]
}%
כלומר \L{$p$} מאפס את \L{$T|_{U}$}: \L{%
\[
p\left(T|_{U}\right)=0_{\mathcal{L}(U)}
\]
}%

יהי \L{$\widetilde{q}_{\min}$} הפולינום המינימלי של \L{$T|_{U}$}.

לפי המשפט מהרצאה \L{$5$}, הפולינום המינימלי מחלק כל פולינום שמאפס
את האופרטור.

לכן \L{$\widetilde{q}_{\min}$} מחלק את \L{%
\[
p(z)=(z-\lambda_{1})\cdots(z-\lambda_{m-1})
\]
}%
מכאן \L{%
\[
\widetilde{q}_{\min}(z)=(z-\alpha_{1})\cdots(z-\alpha_{k})
\]
}%
כאשר \L{$\alpha_{1},\dots,\alpha_{k}$} שונים זה מזה וגם \L{%
\[
\alpha_{1},\dots,\alpha_{k}\in\{\lambda_{1},\dots,\lambda_{m-1}\}
\]
}%
כאשר \L{%
\[
1\leq k\leq m-1
\]
}%
\L{$k\geq1$} כי הפולינום המתוקן היחיד ממעלה \L{$0$} הוא הקבוע \L{$1$},
והוא לא מאפס את \L{$T|_{U}$}: הצבה בו נותנת את אופרטור הזהות על \L{$U$},
שאינו אופרטור האפס כי \L{$\dim U\geq1$}.

מ-)א(, )ב( ו-)ג( נובע ש-\L{$U$} ו-\L{$T|_{U}$} מקיימים את התנאים
של הנחת האינדוקציה )שלב 2(, כלומר היא חלה עבור \L{%
\[
\widetilde{V}=U,\quad\widetilde{T}=T|_{U}
\]
}%

4. לכן קיים בסיס \L{$\mathcal{B}_{U}$} ב-\L{$U$} המורכב מוקטורים
עצמיים של \L{$T$}: \L{%
\[
\mathcal{B}_{U}=(u_{1},\dots,u_{p})
\]
}%
כאשר \L{$u_{1},\dots,u_{p}$} הם וקטורים עצמיים של \L{$T$}.

הנחת האינדוקציה נותנת ב-\L{$U$} בסיס של וקטורים עצמיים של \L{$T|_{U}$}.

אלה גם וקטורים עצמיים של \L{$T$}, כי בתוך \L{$U$} לא משנה אם מדברים
על \L{$T$} או על הצמצום.

עכשיו נרצה להרחיב את זה לבסיס של המרחב המקורי.

5. נסתכל על האופרטור \L{$T-\lambda_{m}I\in\mathcal{L}(V)$}.

לפי משפט מלינארית \L{$1$}, המימד של התחום של כל ה\textquotedblright ל
שווה למימד הגרעין שלה ועוד מימד התמונה שלה: \L{%
\[
\dim V=\dim\text{Null}(T-\lambda_{m}I)+\dim\text{Range}(T-\lambda_{m}I)
\]
}%
נרצה להוכיח ש- \L{%
\[
V=\text{Null}(T-\lambda_{m}I)\oplus\text{Range}(T-\lambda_{m}I)
\]
}%
כי אז נוכל לצרף את הבסיס של \L{$U=\text{Range}(T-\lambda_{m}I)$}
משלב 4 לבסיס של \L{$\text{Null}(T-\lambda_{m}I)$}, ולקבל בסיס של
\L{$V$} המורכב מוקטורים עצמיים של \L{$T$}.

6. נוכיח ש- \L{%
\[
\text{Null}(T-\lambda_{m}I)\cap U=\{0_{V}\}
\]
}%
יהי \L{$u\in\text{Null}(T-\lambda_{m}I)\cap U$}.

מכיוון ש- \L{$u\in\text{Null}(T-\lambda_{m}I)$}: \L{%
\[
T(u)=\lambda_{m}u
\]
}%
לכן לכל \L{$i\in\{1,\dots,m-1\}$}: \L{%
\[
(T-\lambda_{i}I)(u)=\lambda_{m}u-\lambda_{i}u=(\lambda_{m}-\lambda_{i})u
\]
}%
מכיוון ש-\L{$u\in U$}, לפי שלב 3: \L{%
\[
\left((T-\lambda_{1}I)\cdots(T-\lambda_{m-1}I)\right)(u)=0_{V}
\]
}%
משתי השורות האחרונות נקבל \L{%
\[
0_{V}=(\lambda_{m}-\lambda_{1})\cdots(\lambda_{m}-\lambda_{m-1})u
\]
}%
\L{$\lambda_{1},\dots,\lambda_{m}$} שונים זה מזה, לכן כל הסקלרים
\L{$\lambda_{m}-\lambda_{i}$} שונים מ-\L{$0_{\mathbb{F}}$}, ומכאן
\L{%
\[
u=0_{V}
\]
}%
כנדרש.

{\beginL 7\endL}. קיבלנו שהחיתוך של שני התמ\textquotedblright וים
הוא \L{$\left\{ 0_{V}\right\} $}, ולפי ההגדרה מלינארית {\beginL 1\endL}
זה אומר שהסכום שלהם ישר.

מימד של סכום ישר שווה לסכום המימדים של התמ\textquotedblright וים )משפט
המימדים הראשון ומימד החיתוך הוא אפס(, ולפי {\beginL 5\endL} סכום המימדים
של התמ\textquotedblright וים הוא המימד של \L{$V$}.

קיבלנו שמימד הסכום הישר הוא המימד של \L{$V$}.

הסכום הישר הוא תמ\textquotedblright ו של \L{$V$} מהמימד של \L{$V$},
לכן קיבלנו שהוא שווה ל-\L{$V$}.

\L{%
\[
V=\text{Null}(T-\lambda_{m}I)\oplus\text{Range}(T-\lambda_{m}I)
\]
}%

8. לפי שלב 4, ל- \L{$U=\text{Range}(T-\lambda_{m}I)$} יש בסיס המורכב
מוקטורים עצמיים של \L{$T$}. 

גם ל-\L{$\text{Null}(T-\lambda_{m}I)$} יש בסיס כזה.

הוא תמ\textquotedblright ו של \L{$V$}, ולכן לפי לינארית \L{$1$}
יש לו בסיס.

יהי \L{$w$} וקטור בבסיס הזה.

אז \L{$w\neq0_{V}$}, כי וקטורי בסיס הם בת\textquotedblright ל, ומכיוון
ש- \L{$w\in\text{Null}(T-\lambda_{m}I)$} מתקיים \L{%
\[
T(w)=\lambda_{m}w
\]
}%
כלומר \L{$w$} הוא וקטור עצמי של \L{$T$} עם הערך העצמי \L{$\lambda_{m}$}.

לפי שלב 7 הסכום ישר, ולכן שני הבסיסים יחד הם בסיס של \L{$V$}.

קיבלנו בסיס של \L{$V$} המורכב מוקטורים עצמיים של \L{$T$}. 

לפי המשפט על התנאים השקולים ללכסינות מהרצאה \L{$7$}, \L{$(\text{I})$}
מתקיים, כנדרש.
\end{proof}

\part{פונקציות של מטריצות וסיום הפרק}

\section{דוגמה - חזקות ואקספוננט של מטריצה לכסינה}

השימוש העיקרי של לכסון הוא שהוא מאפשר להגדיר פונקציות של מטריצות,
למשל כדי לפתור מערכות של משוואות דיפרנציאליות.

נניח כי \L{%
\[
A=QDQ^{-1}
\]
}%
כאשר \L{%
\[
D=\begin{pmatrix}\lambda_{1} &  &  & 0_{\mathbb{F}}\\
 & \lambda_{2}\\
 &  & \ddots\\
0_{\mathbb{F}} &  &  & \lambda_{n}
\end{pmatrix}
\]
}%
)הערכים העצמיים לא חייבים להיות שונים זה מזה(.

נחשב למה שווה חזקה כלשהי של \L{$A$}.\L{%
\begin{align*}
A^{p} & =\underset{p\text{ times}}{\underbrace{(QDQ^{-1})(QDQ^{-1})\cdots(QDQ^{-1})}}\\
 & =(QDIDID\dots IDQ^{-1})=QD^{p}Q^{-1}\\
 & =Q\begin{pmatrix}\lambda_{1}^{p} &  & 0_{\mathbb{F}}\\
 & \ddots\\
0_{\mathbb{F}} &  & \lambda_{n}^{p}
\end{pmatrix}Q^{-1}
\end{align*}
}%

נזכיר מאינפי {\beginL 2\endL} את טור טיילור של \L{$e^{x}$} סביב \L{$0$}:
\L{%
\[
e^{x}=\sum_{p=0}^{\infty}\frac{x^{p}}{p!}
\]
}%
הטור הזה מתכנס לכל \L{$x$}.

באגף ימין יש רק חזקות, חילוק בסקלר וסכום, ואלה פעולות שמוגדרות גם
עבור מטריצות.

לכן נגדיר את \L{$e^{A}$} על ידי הצבת \L{$A$} במקום \L{$x$}. 

נקבל שלפי החישוב של \L{$A^{p}$}: \L{%
\[
\sum_{p=0}^{\infty}\frac{A^{p}}{p!}=\sum_{p=0}^{\infty}\frac{QD^{p}Q^{-1}}{p!}=Q\left(\sum_{p=0}^{\infty}\frac{D^{p}}{p!}\right)Q^{-1}
\]
}%
כל \L{$D^{p}$} אלכסונית, ולכן הסכום שבסוגריים מחושב איבר-איבר על
האלכסון: \L{%
\begin{align*}
e^{A} & \overset{\text{def}}{=}Q\left(\sum_{p=0}^{\infty}\frac{D^{p}}{p!}\right)Q^{-1}\\
 & =Q\begin{pmatrix}\underset{p=0}{\overset{\infty}{\sum}}\frac{\lambda_{1}^{p}}{p!} &  & 0_{\mathbb{F}}\\
 & \ddots\\
0_{\mathbb{F}} &  & \underset{p=0}{\overset{\infty}{\sum}}\frac{\lambda_{n}^{p}}{p!}
\end{pmatrix}Q^{-1}\\
 & =Q\begin{pmatrix}e^{\lambda_{1}} &  & 0_{\mathbb{F}}\\
 & \ddots\\
0_{\mathbb{F}} &  & e^{\lambda_{n}}
\end{pmatrix}Q^{-1}
\end{align*}
}%
כך אפשר להגדיר כל פונקציה שמיוצגת על ידי טור, כל עוד הטור מתכנס בכל
ערך עצמי של \L{$A$}.

\section{הערה - סיום הפרק והמבחן}

בזה הסתיים הפרק על לכסון מטריצות.

בשבוע הבא מתחילים מרחבי מכפלה פנימית.\\

על המבחן: לכל פרק תפורסם רשימת משפטים, ומהפרק הזה נכנסו אליה חמישה.

בחלק התיאורטי יופיעו שלושה משפטים מהרשימה, בניסוח של המרצה, וצריך
לבחור שניים מהם.

כל שאלה שווה \L{$25$} נקודות )למשל הוכחת משפט \L{$20$} נקודות ועוד
הגדרה \L{$5$} נקודות(, ובסך הכל \L{$50$} נקודות.

\part{נספח}

\section{הוכחה ישירה של הכיוון הראשון במשפט - לכסינות באמצעות ריבוי אלגברי
וריבוי גאומטרי}

יהיו \L{$\lambda_{1},\dots,\lambda_{m}$} כל הערכים העצמיים של \L{$T$}
)שונים זה מזה(, כלומר\L{%
\[
\forall i,j\in\{1,\dots,m\}\text{ s.t. }i\neq j:\quad\lambda_{j}\neq\lambda_{i}
\]
}%

ונניח כי קיים בסיס \L{$\mathcal{B}_{V}$} ב-\L{$V$} כך ש- \L{%
\[
M(T,\mathcal{B}_{V},\mathcal{B}_{V})=\begin{pmatrix}a_{1,1} & 0_{\mathbb{F}} & \cdots & 0_{\mathbb{F}}\\
0_{\mathbb{F}} & a_{2,2} &  & \vdots\\
\vdots &  & \ddots\\
0_{\mathbb{F}} & 0_{\mathbb{F}} &  & a_{n,n}
\end{pmatrix}
\]
}%

נוכיח ש:

)א( סכום הריבויים האלגבריים של הערכים העצמיים שווה למימד של \L{$V$}:\L{%
\[
a_{\left(\lambda_{1}\right)}+\dots+a_{\left(\lambda_{m}\right)}=n
\]
}%

)ב( עבור כל ערך עצמי, הריבוי האלגברי שלו שווה לריבוי הגיאומטרי שלו:

\L{%
\[
g_{\left(\lambda_{1}\right)}=a_{\left(\lambda_{1}\right)},\dots,g_{\left(\lambda_{m}\right)}=a_{\left(\lambda_{m}\right)}
\]
}%

\begin{proof}
1. לפי המשפט על התנאים השקולים ללכסינות מהרצאה {\beginL 7\endL}, קיים
בסיס \L{$\mathcal{B}_{V}$} המורכב מוקטורים עצמיים של \L{$T$}.

2. נסדר את הוקטורים מחדש לפי הערכים העצמיים שלהם )כמו בהרצאה {\beginL 7\endL}(.

לכל \L{$i\in\{1,\dots,m\}$}, נסמן את וקטורי הבסיס השייכים ל- \L{$\lambda_{i}$}
כך: \L{%
\[
\mathcal{B}^{(i)}=(v_{1}^{(i)},\dots,v_{n_{i}}^{(i)})
\]
}%
האינדקס העליון אומר לאיזה ערך עצמי הוקטור שייך, והתחתון ממספר את הוקטורים
בתוך הקבוצה. 

נשים את הקבוצות אחת אחרי השנייה: \L{%
\[
\mathcal{B}'_{V}=(\mathcal{B}^{(1)},\mathcal{B}^{(2)},\dots,\mathcal{B}^{(m)})
\]
}%
\L{$\mathcal{B}'_{V}$} הוא בסיס של \L{$V$}, כי אלה אותם וקטורים
בסדר אחר.

לא החסרנו ולא הוספנו וקטורים מהבסיס לכן\L{%
\[
n_{1}+\dots+n_{m}=n
\]
}%

3. מאיך שהגדרנו אותו, לכל וקטור בבסיס \L{$v_{j}^{(i)}$} מתקיים \L{$T\left(v_{j}^{(i)}\right)=\lambda_{i}v_{j}^{(i)}$}.

אז מהגדרת המטריצה המייצגת \L{$\lambda_{1}$} מופיע \L{$n_{1}$} פעמים,
\L{$\lambda_{2}$} מופיע \L{$n_{2}$} פעמים וכן הלאה עד \L{$\lambda_{m}$}
שמופיע \L{$n_{m}$} פעמים, ובכל מקום אחר יש אפס.

\L{%
\[
M(T,\mathcal{B}'_{V},\mathcal{B}'_{V})=\begin{pmatrix}\lambda_{1} &  &  &  &  &  & 0_{\mathbb{F}}\\
 & \ddots\\
 &  & \lambda_{1}\\
 &  &  & \ddots\\
 &  &  &  & \lambda_{m}\\
 &  &  &  &  & \ddots\\
0_{\mathbb{F}} &  &  &  &  &  & \lambda_{m}
\end{pmatrix}
\]
}%

מכאן, \L{%
\[
P_{ch}(z)=\det\left(zI_{n}-M(T,\mathcal{B}'_{V},\mathcal{B}'_{V})\right)=(z-\lambda_{1})^{n_{1}}\cdots(z-\lambda_{m})^{n_{m}}
\]
}%
\L{$\lambda_{1},\dots,\lambda_{m}$} שונים זה מזה, לכן החזקה הגדולה
ביותר של \L{$(z-\lambda_{i})$} שמחלקת את \L{$P_{ch}$} היא בדיוק
\L{$n_{i}$} ונקבל

\L{%
\[
a_{\left(\lambda_{1}\right)}=n_{1},\dots,a_{\left(\lambda_{m}\right)}=n_{m}
\]
}%
\L{%
\[
\therefore\quad a_{\left(\lambda_{1}\right)}+\dots+a_{\left(\lambda_{m}\right)}=n_{1}+\dots+n_{m}=n
\]
}%

שזה )א(.

{\beginL 4\endL}. כל \L{$v_{j}^{(i)}$} הוא וקטור עצמי של \L{$\lambda_{i}$},
כלומר שייך ל- \L{$\text{Null}(T-\lambda_{i}I)$}, ו- \L{$\text{Null}(T-\lambda_{i}I)$}
הוא תמ\textquotedblright ו ולכן מכיל גם את כל הצירופים הלינאריים שלהם.\L{%
\begin{align*}
\text{Span}\left(v_{1}^{(1)},\dots,v_{n_{1}}^{(1)}\right) & \subseteq\text{Null}(T-\lambda_{1}I)\\
 & \vdots\\
\text{Span}\left(v_{1}^{(m)},\dots,v_{n_{m}}^{(m)}\right) & \subseteq\text{Null}(T-\lambda_{m}I)
\end{align*}
}%

{\beginL 5\endL}. \L{$\mathcal{B}'_{V}$} הוא בסיס של \L{$V$} שמחולק
לקבוצות \L{$\mathcal{B}^{(1)},\dots,\mathcal{B}^{(m)}$}, להן אין
וקטורים משותפים.

כל \L{$v\in V$} נכתב באופן יחיד כצירוף לינארי של וקטורי הבסיס.

אם מקבצים את המחוברים לפי הקבוצות מקבלים הצגה יחידה \L{$v=u_{1}+\dots+u_{m}$}
כאשר \L{$u_{i}\in\text{Span}\left(\mathcal{B}^{(i)}\right)$}.

לכן \L{$V$} הוא סכום ישר של ה-\L{$\text{Span}$}-ים )זו בדיוק ההגדרה
של אקסלר לסכום ישר(.

\L{%
\[
V=\text{Span}\left(v_{1}^{(1)},\dots,v_{n_{1}}^{(1)}\right)\oplus\dots\oplus\text{Span}\left(v_{1}^{(m)},\dots,v_{n_{m}}^{(m)}\right)
\]
}%

{\beginL 6\endL}. לפי המשפט על התנאים השקולים ללכסינות מהרצאה \L{$7$},
\L{%
\[
V=\text{Null}(T-\lambda_{1}I)\oplus\dots\oplus\text{Null}(T-\lambda_{m}I)
\]
}%

7. יהי \L{$i\in\{1,\dots,m\}$}. נוכיח שמתקיים \L{%
\[
\text{Span}\left(v_{1}^{(i)},\dots,v_{n_{i}}^{(i)}\right)=\text{Null}(T-\lambda_{i}I)
\]
}%
בעזרת הכלה דו-כיוונית.

\L{$\text{Span}\left(\mathcal{B}^{(i)}\right)\subseteq\text{Null}(T-\lambda_{i}I)$}:
הוכחנו ב-){\beginL 4\endL}(.

\L{$\text{Null}(T-\lambda_{i}I)\subseteq\text{Span}\left(\mathcal{B}^{(i)}\right)$}:

יהי \L{$u\in\text{Null}(T-\lambda_{i}I)$}.

לפי ){\beginL 5\endL}(, קיימים \L{$u_{1},\dots,u_{m}$} כך ש-

\L{%
\[
u=\underset{\in\text{Span}\left(v_{1}^{(1)},\dots,v_{n_{1}}^{(1)}\right)}{\underbrace{u_{1}}}+\quad\dots\quad+\underset{\in\text{Span}\left(v_{1}^{(m)},\dots,v_{n_{m}}^{(m)}\right)}{\underbrace{u_{m}}}
\]
}%

ולפי ){\beginL 4\endL}(, לכל \L{$k\in\left\{ 1,\dots,m\right\} $}
מתקיים \L{$u_{k}\in\text{Null}(T-\lambda_{k}I)$}.

מצד שני, \L{$u\in\text{Null}(T-\lambda_{i}I)$} ו-\L{$0_{V}$} שייך
לכל תמ\textquotedblright ו, לכן אפשר לכתוב את \L{$u$} גם כך )\L{$u$}
במקום ה-\L{$i$}, ו-\L{$0_{V}$} בכל שאר המקומות(: \L{%
\[
u=\underbrace{0_{V}}_{\in\text{Null}(T-\lambda_{1}I)}+\dots+\underbrace{u}_{\in\text{Null}(T-\lambda_{i}I)}+\dots+\underbrace{0_{V}}_{\in\text{Null}(T-\lambda_{m}I)}
\]
}%

קיבלנו שתי הצגות של \L{$u$} כסכום של וקטורים מ- \L{$\text{Null}(T-\lambda_{1}I),\dots,\text{Null}(T-\lambda_{m}I)$}.

לפי ){\beginL 6\endL}( הסכום של תתי-המרחבים האלה ישר, ולכן ההצגה יחידה,
כלומר שתי ההצגות שוות רכיב-רכיב.

בפרט, ברכיב ה-\L{$i$}: \L{%
\[
u_{i}=u
\]
}%
לכן \L{%
\[
u=u_{i}\in\text{Span}\left(v_{1}^{(i)},\dots,v_{n_{i}}^{(i)}\right)
\]
}%
כנדרש.

קיבלנו: \L{%
\begin{align*}
\text{Span}\left(v_{1}^{(1)},\dots,v_{n_{1}}^{(1)}\right) & =\text{Null}(T-\lambda_{1}I)\\
 & \vdots\\
\text{Span}\left(v_{1}^{(m)},\dots,v_{n_{m}}^{(m)}\right) & =\text{Null}(T-\lambda_{m}I)
\end{align*}
}%

{\beginL 8\endL}. הוקטורים של \L{$\mathcal{B}^{(i)}$} הם חלק מהבסיס
\L{$\mathcal{B}'_{V}$} ולכן בת\textquotedblright ל.

לכן הרשימה \L{$\mathcal{B}^{(i)}=\left(v_{1}^{(i)},\dots,v_{n_{i}}^{(i)}\right)$}
היא בת\textquotedblright ל ופורשת, דהיינו בסיס של \L{$\text{Span}\left(\mathcal{B}^{(i)}\right)$},
ומתקיים\L{%
\[
\dim\left(\text{Span}\left(\mathcal{B}^{(i)}\right)\right)=n_{i}
\]
}%

לפי ){\beginL 7\endL}(, מתקיים \L{$\dim\left(\text{Null}(T-\lambda_{i}I)\right)=\dim\left(\text{Span}\left(\mathcal{B}^{(i)}\right)\right)=n_{i}$}

9. לפי ){\beginL 3\endL}( מתקיים \L{$a_{(\lambda_{i})}=n_{i}$}, ולפי
ההגדרה של הריבוי הגאומטרי, \L{$g_{(\lambda_{i})}=\dim\left(\text{Null}(T-\lambda_{i}I)\right)$}.

אז בסה\textquotedblright כ\L{%
\[
\forall1\le i\le m:
\]
}%
\L{%
\[
a_{(\lambda_{i})}=n_{i}=\dim\left(\text{Null}(T-\lambda_{i}I)\right)=g_{(\lambda_{i})}
\]
}%
שזה )ב(, כנדרש.
\end{proof}

\end{document}
